%----------
%	ANALISIS GENERAL - MACHINE LEARNING
%----------	

A total of 163 experiments have been conducted for \textit{"Análisis General}" in the span of five months (12/02/24 - 08/07/24), defining ten different windows that can be seen in Figures ~\ref{fig:exp_analisis_general} and ~\ref{fig:exp_analisis_general_classes} and will be further described below. In particular, Figure \ref{fig:exp_analisis_general} shows the overall trend in weighted average F1 scores across different experimental windows, marked by significant changes. The linear trend line, with $R^2 = 0.07$, underscores the moderate yet positive impact of these incremental improvements over time. Figure \ref{fig:exp_analisis_general_classes} provides a detailed class-specific breakdown, highlighting the individual performances for `Comentario Positivo' and `Comentario Negativo'. While both classes exhibit a positive trend, the metrics for `Comentario Positivo' are notably more stable. However, an exception is observed in W3, where the performance decreases significantly, marking the worst outcome of the experimentation.

These insights can also be seen in Table~\ref{tab:window-metrics-ageneral}, where we notice W10 as the window with the best mean performance metrics, with weighted average F1-Score of 0.931, and 0.955 and 0.847 for 'Comentario Positivo' and 'Comentario Negativo', respectively.

\begin{figure}[H]
	\ffigbox[\FBwidth]
	{\caption{Distribution of the ML Experiments over Time for "Análisis General" with the Different Windows and Weighted Avg. F1-Score Trend Line}\label{fig:exp_analisis_general}}
 % [Aquí ponemos como queremos que aparezca en el índice, sin la referencia]{Aquí como queremos que aparezca en el documento, con la referencia}
	{\includegraphics[scale=0.25]{Plantilla_TFG_ingles_2019/imagenes/Experiments in _Análisis General_ over time.pdf}}
\end{figure}

\begin{figure}[H]
	\ffigbox[\FBwidth]
	{\caption{Distribution of the ML Experiments over Time for "Análisis General", per Class and with the Different Windows (Trend Line of W. Avg. F1-Score in Light Red)}\label{fig:exp_analisis_general_classes}}
 % [Aquí ponemos como queremos que aparezca en el índice, sin la referencia]{Aquí como queremos que aparezca en el documento, con la referencia}
	{\includegraphics[scale=0.25]{Plantilla_TFG_ingles_2019/imagenes/Experiments in _Análisis General_ over time_classes_trendline2.pdf}}
\end{figure}


\begin{table}[H]
    \centering
    \begingroup
    \small % Change to desired font size
    \ttabbox[\FBwidth]{
        \caption{AG Average Evaluation Metrics per Window in ML Experiments}
        \label{tab:window-metrics-ageneral} % Unique label for the table
    }{
        \begin{tabularx}{\textwidth}{
            X
            >{\centering\arraybackslash}X
            >{\centering\arraybackslash}X
            >{\centering\arraybackslash}X
            >{\centering\arraybackslash}X
            >{\centering\arraybackslash}X
            >{\centering\arraybackslash}X
            >{\columncolor{gray!15}\centering\arraybackslash}X % Highlight this column
        }
            \hline
            \rowcolor{gray!30} % Adding a gray color to the header row
            \textbf{} & \textbf{Acc. (Std.)} & \textbf{Time (s)} & \textbf{W. Avg. Precision} & \textbf{W. Avg. Recall} & \textbf{W. Avg. F1-Score} & \textbf{C. Pos. F1-Score} & \textbf{C. Neg. F1-Score} \\
            \hline
            W1 & 0.795 (0.012) & 2303.864 & 0.815 & 0.795 & 0.802 & 0.879 & 0.599 \\
            \hline
            W2 & 0.834 (0.019) & 1094.649 & 0.832 & 0.834 & 0.832 & 0.904 & 0.592 \\
            \hline
            W3 & 0.713 (0.038) & 479.563 & 0.806 & 0.713 & 0.718 & 0.778 & 0.428 \\
            \hline
            W4 & 0.819 (0.029) & 1102.832 & 0.825 & 0.819 & 0.818 & 0.882 & 0.594 \\
            \hline
            W5 & 0.820 (0.025) & 1062.982 & 0.828 & 0.820 & 0.819 & 0.883 & 0.597 \\
            \hline
            W6 & 0.822 (0.018) & 799.329 & 0.834 & 0.822 & 0.818 & 0.884 & 0.585 \\
            \hline
            W7 & 0.808 (0.012) & 898.133 & 0.845 & 0.808 & 0.812 & 0.861 & 0.640 \\
            \hline
            W8 & 0.814 (0.027) & 955.299 & 0.853 & 0.814 & 0.816 & 0.864 & 0.645 \\
            \hline
            W9 & 0.813 (0.024) & 983.487 & 0.855 & 0.813 & 0.817 & 0.862 & 0.655 \\
            \hline
            W10 & 0.931 (0.009) & 6776.630 & 0.933 & 0.931 & 0.931 & 0.955 & 0.847 \\
            \hline
            \multicolumn{8}{l}{Source: 'Training.xlsx'} \\
            \hline
        \end{tabularx}
    }
    \endgroup
\end{table}


\newpage

\begin{itemize}
    \item \textbf{W1: Baseline Establishment:} Implementation of two fundamental ML models, including Random Forest and Logistic Regression, utilizing traditional embeddings such as FastText and Bag of Words (BoW). This phase also involved the combination of no balancing technique with upsampling to establish a robust baseline for further experimentation. From Table~\ref{tab:window-metrics-ageneral} we notice a weighted average F1-Score of 0.802, with individual F1-Scores of 0.879 for "Comentario Positivo" and 0.599 for "Comentario Negativo", indicating a decent start but with room for improvement, particularly for negative comments.

    
    \item \textbf{W2-W3: Introduction of Advanced Embeddings and Preprocessing:}
        \begin{itemize}
            \item \textbf{W2:} Implemented the 'RoBERTa' embedding described in Section~\ref{text_encoding_dl}, and included additional stylistic features detailed in Section~\ref{stylistic_features} like filtering tweets in Spanish, adding the tweet length, and the number of adjectives. These modifications improved the overall F1 score (0.802 vs 0.832) but did not significantly enhance metrics for 'Comentario Negativo' (0.599 vs 0.592). 

            \item \textbf{W3:} Implemented more traditional ML models such as XGBoosting, MLP, and naive Bayes, along with the "RoBERTuito" embedding described in Section~\ref{text_encoding_dl}. Despite these additions, the overall performance declined further to a weighted average F1-Score of 0.718, with an F1-Score of 0.428 for 'Comentario Negativo', marking this as the worst result of the experimentation phase.
        \end{itemize}

    \item \textbf{W4-W9: Text Pre-Processing and Lexicons:}
        \begin{itemize}
            \item \textbf{W4:} Applied comprehensive text preprocessing as explained in Section~\ref{sec:data_preprocessing}, including normalizing text to lowercase, removing emojis, numbers, and punctuation, eliminating non-contributive words, reducing elongated words and repeated characters. This window maintained a similar performance to W2, with a weighted average F1-Score of 0.818, and F1-Scores of 0.882 for 'Comentario Positivo' and 0.594 for 'Comentario Negativo'.
        
            \item \textbf{W5:} Introduced the lexicons detailed in Section~\ref{lexicons} with the ADASYN balancing technique for all ML models. The performance metrics for both classes showed slight improvements, with a weighted average F1-Score of 0.819, and balanced F1-Scores for 'Comentario Positivo' (0.883) and 'Comentario Negativo' (0.597).
        
            \item \textbf{W6:} Conducted experiments for all models with no balancing technique to evaluate the effects of the previous modifications. The performance for 'Comentario Negativo' decreased slightly to an F1-Score of 0.585, with a weighted average F1-Score of 0.818.
        
            \item \textbf{W7:} Experimented with upsampling and traditional embeddings for all models, where we noticed an improvement for 'Comentario Negativo' F1-Score (0.640), but lower performance for "Comentario Positivo" and weighted average F1-Score, dropping to 0.861 and 0.812, respectively.
        
            \item \textbf{W8:} Continued with similar experimentation as W7, but applying SMOTE. It achieved a slight recovery in performance metrics compared to W7, with the weighted average F1-Score increasing to 0.816, 0.864 for 'Comentario Positivo' and 0.645 for 'Comentario Negativo'.
        
            \item \textbf{W9:} Similar settings as in W7 and W8 but with ADASYN technique. The results were stable for the weighted average F1-Score (0.816 vs 0.817) and "Comentario Positivo" (0.864 vs 0.862), but better for 'Comentario Negativo' (0.645 vs 0.655).
        \end{itemize}
        
    \item \textbf{W10: Introduced OpenAI Embedding:} From Table~\ref{tab:window-metrics-ageneral}, the metrics and F1 scores showed superior performance for both classes. This window achieved a significant increase in F1 scores, reaching 0.955 for 'Comentario Positivo' and 0.847 for 'Comentario Negativo', resulting in the best overall performance with a weighted average F1-Score of 0.931. These results suggest the high quality of this embedding in capturing the particular context of the dataset, but also the high computational time it required (6776.630 seconds).
        
\end{itemize}



Taking everything into account, we provide the following detailed tables with the top 5 best overall models for \textit{Análisis General} (AG). Table \ref{tab:ageneral-bestmodels} lists the models, and their configurations, and Table \ref{tab:ageneral-bestmodels_metrics} presents the performance metrics for these models. Notice that we have included two naive models in both tables in order to compare them to our models. {\agmajor} is a classifier that only predicts the majority class, while {\agrandom} predicts randomly, based on the class probabilities.

%It is important to point out that for these models we are assuming that the Standard Deviation is zero because we are using the probabilities of each class, and that F1-Score is approximately equal to such probabilities.

\begin{table}[H]
    \centering
    \begingroup
    \small % Change to desired font size
    \ttabbox[\FBwidth]{
        \caption{AG Best ML Models Nomenclature}
        \label{tab:ageneral-bestmodels} % Unique label for the table
    }{
        \begin{tabularx}{\textwidth}{
            X  % Name
            X  % Model
            X  % Embedding
            >{\centering\arraybackslash}X  % CV
            X  % Balance
            }
            \hline
            \rowcolor{gray!30} % Adding a gray color to the header row
            \textbf{} & \textbf{Model} & \textbf{Embedding} & \textbf{CV} & \textbf{Balance} \\
            \hline
            \multicolumn{5}{l}{\textbf{Naive Models}} \\
            \hline 
            \agmajor & - & - & - & - \\
            \hline
            \agrandom & - & - & - & - \\
            \hline
            \multicolumn{5}{l}{\textbf{Finetuned Models}} \\
            \hline
            \agmlpteup & MLP & text-embedding-3-large & 5 & Upsampling \\
            \hline
            \agrftesm & Random Forest & text-embedding-3-large & 5 & SMOTE \\
            \hline
            \aglrteno & Logistic Regression & text-embedding-3-large & 5 & None \\
            \hline
            \agxgbtesm & XGBoosting & text-embedding-3-large & 5 & SMOTE \\
            \hline
            \aglrteup & Logistic Regression & text-embedding-3-large & 5 & Upsampling \\
            \hline
            \multicolumn{5}{l}{Source: 'Training.xlsx'} \\
            \hline
        \end{tabularx}
    }
    \endgroup
\end{table}




\begin{table}[H]
    \centering
    \begingroup
    \small % Change to desired font size
    \ttabbox[\FBwidth]{
        \caption{AG Evaluation Metrics for the Best ML Models}
        \label{tab:ageneral-bestmodels_metrics} % Unique label for the table
    }{
        \begin{tabularx}{\textwidth}{
            X
            >{\centering\arraybackslash}X
            >{\centering\arraybackslash}X
            >{\centering\arraybackslash}X
            >{\columncolor{gray!15}\centering\arraybackslash}X % Highlight this column
        }
            \hline
            \rowcolor{gray!30} % Adding a gray color to the header row
            \textbf{} & \textbf{Acc. (Std.)} & \textbf{Weighted Avg. F1-Score} & \textbf{C. Positivo F1-Score} & \textbf{C. Negativo F1-Score} \\
            \hline
            \multicolumn{5}{l}{\textbf{Naive Models}} \\
            \hline
            \agmajor & 0.778 (0.010) & 0.680 & 0.875 & 0.000 \\
            \hline
            \agrandom & 0.654 (0.012) & 0.654 & 0.778 & 0.222 \\
            \hline
            \multicolumn{5}{l}{\textbf{Finetuned Models}} \\
            \hline
            \agmlpteup & 0.945 (0.007) & 0.946 & 0.965 & 0.878 \\
            \hline
            \agrftesm & 0.942 (0.008) & 0.942 & 0.963 & 0.867 \\
            \hline
            \aglrteno & 0.939 (0.008) & 0.940 & 0.961 & 0.865 \\
            \hline
            \agxgbtesm & 0.939 (0.005) & 0.940 & 0.961 & 0.865 \\
            \hline
            \agxgbtesm & 0.939 (0.007) & 0.939 & 0.961 & 0.863 \\
            \hline
            \multicolumn{5}{l}{Source: 'Training.xlsx'} \\
            \hline
        \end{tabularx}
    }
    \endgroup
\end{table}


We have chosen model \textbf{\agmlpteup}  as the best performing model because it achieved the highest F1-Score for negative comments ("Comentario Negativo") while maintaining strong performance across other metrics, and outperforming naive models {\agmajor} and {\agrandom}. Ensuring a high "Comentario Negativo" F1-Score is crucial, as the second and third classification tasks, \textit{"Contenido Negativo"} and \textit{"Insults"}, focus exclusively on identifying the type of negative comments and insults, respectively.

Figure \ref{fig:exp_analisis_general_metrics} presents the confusion matrix for this model. Notably, it shows a higher number of true negatives (65) compared to false negatives (10), with only 8 false positives. The results suggest that the combination of upsampling techniques and advanced embeddings (\textit{"text-embedding-3-large"}) significantly enhanced the model's performance, making this ML approach highly successful for this classification task.


\begin{figure}[H]
	\ffigbox[\FBwidth]
	{\caption{Confusion Matrix for the Best ML Model (\textbf{\agmlpteup}) for "Análisis General"}\label{fig:exp_analisis_general_metrics}}
 % [Aquí ponemos como queremos que aparezca en el índice, sin la referencia]{Aquí como queremos que aparezca en el documento, con la referencia}
	{\includegraphics[scale=0.14]{Plantilla_TFG_ingles_2019/imagenes/Experiments in _Análisis General_ best_model.pdf}}
\end{figure}



